\documentclass[12pt]{article}
\usepackage[utf8]{inputenc}
\usepackage[T1]{fontenc}
\usepackage[margin=1in]{geometry}
\usepackage{setspace}
\usepackage{xcolor}
\usepackage[colorlinks=true, urlcolor=blue!60!black, linkcolor=blue!60!black]{hyperref}
\setlength{\parindent}{0em}
\setlength{\parskip}{0.45em}
\renewcommand{\baselinestretch}{1.0}
\begin{document}
\begin{center}
{\Large \textit{Garcetti v.\ Ceballos}, 547 U.S.\ 410 (2006)\par}
\vspace{0.3em}
{\normalsize Case Brief --- One-Page Edition\par}
\end{center}

\textbf{Memory Jogger.} A prosecutor's internal memo was official-duty speech, so the First Amendment did not shield him from discipline (\href{https://supreme.justia.com/cases/federal/us/547/410/}{Justia}).

\textbf{Facts.} Deputy District Attorney Richard Ceballos, a calendar deputy in Los Angeles County, wrote a memo recommending dismissal of a case because the search-warrant affidavit contained serious misrepresentations. He alleged his supervisors retaliated with a reassignment, a transfer, and a denied promotion.

\textbf{Procedural History.} The district court granted summary judgment for the supervisors; the Ninth Circuit reversed.

\textbf{Votes.} Reversed 5--4: Kennedy for the Court; Stevens, Souter (with Stevens and Ginsburg), and Breyer dissented.

\textbf{Holding.} Statements made pursuant to official duties are not citizen speech and are unprotected; \textit{Pickering} balancing never begins. Scholarship and teaching were reserved.

\textbf{Reasoning.} Such speech is work product the employer ``commissioned or created,'' and protecting it would invite judicial supervision of agencies; whistleblower laws offer other remedies. Stevens answered ``Sometimes,'' not ``Never''; Souter would protect speech on matters of ``unusual importance.''

\textbf{Examples.} \emph{Same-Side:} an inspector's internal fire-code report. \emph{Opposite-Side:} subpoenaed trial testimony (\textit{Lane v.\ Franks}). \emph{Fence-Sitter:} a professor's scholarship.

\textbf{Critique.} The rule punishes employees who report wrongdoing up the chain while leakers may keep protection, leaving whistleblowers to patchy statutes; supporters say it keeps judges from managing government workplaces.

\textbf{Key Quotation.} ``We hold that when public employees make statements pursuant to their official duties, the employees are not speaking as citizens for First Amendment purposes, and the Constitution does not insulate their communications from employer discipline.'' --- Justice Kennedy, 547 U.S.\ at 421.

\end{document}
